\documentclass[11pt]{article}
\usepackage{amsmath,amssymb,amsthm}
\usepackage{bm}
\usepackage{algorithm}
\usepackage{algorithmic}
\usepackage{enumitem}
\usepackage{hyperref}
\usepackage{geometry}
\geometry{margin=1in}
\newtheorem{theorem}{Theorem}
\newtheorem{lemma}{Lemma}
\newtheorem{assumption}{Assumption}
\newtheorem{corollary}{Corollary}
\newcommand{\E}{\mathbb{E}}
\newcommand{\R}{\mathbb{R}}
\newcommand{\norm}[1]{\left\|#1\right\|}
\newcommand{\inner}[2]{\left\langle #1,#2 \right\rangle}
\newcommand{\argmin}{\operatorname*{arg\,min}}

\begin{document}

\section*{Algorithm Name}
\textbf{Randomized Block Coordinate Descent (RBCD)}  

The method repeatedly picks a single block of coordinates at random, computes the (partial) gradient on that block, and takes a step only in the selected directions while leaving all other coordinates unchanged.

\section*{Mathematical Formulation}
\begin{itemize}
    \item Parameter vector $w\in\R^{d}$ is partitioned into $n$ (non‑overlapping) blocks
          \[
          w = \bigl(w^{(1)},w^{(2)},\dots ,w^{(n)}\bigr),\qquad 
          w^{(i)}\in\R^{d_i},\;\; \sum_{i=1}^{n}d_i = d .
          \]
    \item At iteration $t$, draw a block index $i_t\in\{1,\dots ,n\}$ according to a probability
          distribution $p=(p_1,\dots ,p_n)$, $p_i>0$, $\sum_i p_i=1$.
          In the most common case $p_i=1/n$ (uniform sampling).
    \item Define the \emph{block selector} matrix $U_i\in\R^{d\times d_i}$ that extracts block $i$:
          $U_i^{\top}w = w^{(i)}$ and $U_iU_i^{\top}$ projects onto the coordinates of block $i$.
    \item Let $\nabla_i f(w) := U_i^{\top}\nabla f(w)$ be the partial gradient w.r.t. block $i$.
    \item Learning‑rate (step‑size) sequence $\{\eta_t\}_{t\ge 0}$, $\eta_t>0$.
    \item \textbf{Update rule}
          \[
          w_{t+1}= w_t - \eta_t\,U_{i_t}\,\nabla_i f(w_t) .
          \tag{1}
          \]
          Equivalently, for the selected block $i_t$,
          \[
          w_{t+1}^{(i_t)} = w_t^{(i_t)} - \eta_t \,\nabla_{i_t} f(w_t),\qquad 
          w_{t+1}^{(j)} = w_t^{(j)}\;\;(j\neq i_t).
          \]
\end{itemize}

\section*{Pseudocode}
\begin{algorithm}[H]
\caption{Randomized Block Coordinate Descent}
\begin{algorithmic}[1]
\STATE \textbf{Input:} initial point $w_0\in\R^{d}$, step‑size schedule $\{\eta_t\}$, block distribution $p$.
\FOR{$t=0,1,2,\dots,T-1$}
    \STATE Sample block $i_t\sim p$.
    \STATE Compute partial gradient $\displaystyle g_t = \nabla_{i_t} f(w_t)$.
    \STATE Update selected block:
    \[
    w_{t+1}^{(i_t)} = w_t^{(i_t)} - \eta_t\, g_t .
    \]
    \STATE Keep other blocks unchanged:
    \[
    w_{t+1}^{(j)} = w_t^{(j)}\quad \forall j\neq i_t .
    \]
\ENDFOR
\STATE \textbf{Output:} $w_T$ (or the averaged iterate $\bar w_T = \frac1T\sum_{t=0}^{T-1} w_t$).
\end{algorithmic}
\end{algorithm}

\section*{Dry Run Example}
Consider the one‑dimensional quadratic
\[
f(w)= (w-3)^2 ,\qquad w\in\R .
\]
There is only one block ($n=1$), thus $i_t=1$ deterministically.
Take $w_0=0$ and a constant step $\eta_t=\eta=0.1$.

\begin{center}
\begin{tabular}{c|c|c|c}
$t$ & $w_t$ & $\nabla f(w_t)=2(w_t-3)$ & $f(w_t)$\\ \hline
0 & $0$ & $-6$ & $9$\\
1 & $w_1 = 0-0.1(-6)=0.6$ & $2(0.6-3)=-4.8$ & $(0.6-3)^2=5.76$\\
2 & $w_2 = 0.6-0.1(-4.8)=1.08$ & $2(1.08-3)=-3.84$ & $(1.08-3)^2=3.6864$\\
3 & $w_3 = 1.08-0.1(-3.84)=1.464$ & $2(1.464-3)=-3.072$ & $(1.464-3)^2=2.3670$\\
\end{tabular}
\end{center}
We see the objective value decreasing at each iteration, illustrating how the algorithm moves toward the minimiser $w^\star=3$.

\newpage
\section*{Assumptions}
\begin{assumption}[Convexity]\label{as:convex}
$f:\R^{d}\to\R$ is convex, i.e.
\[
f(y)\ge f(x)+\inner{\nabla f(x)}{y-x},\qquad\forall x,y\in\R^{d}.
\]
\end{assumption}

\begin{assumption}[Block‑wise $L_i$‑smoothness]\label{as:smooth}
For each block $i\in\{1,\dots ,n\}$ there exists $L_i>0$ such that
\[
f\bigl(x+U_i s\bigr)\le 
f(x)+\inner{\nabla_i f(x)}{s}+\frac{L_i}{2}\norm{s}^2,
\qquad\forall x\in\R^{d},\; \forall s\in\R^{d_i}.
\tag{2}
\]
\end{assumption}

\begin{assumption}[Bounded distance to optimum]\label{as:bounded}
Let $w^\star\in\arg\min f$ be a minimiser. The initial point satisfies
\[
R^2:=\norm{w_0-w^\star}^2<\infty .
\]
\end{assumption}

\begin{assumption}[Sampling]\label{as:sampling}
Blocks are sampled i.i.d. with probabilities $p_i>0$, $\sum_{i=1}^n p_i=1$.
Define
\[
p_{\min}:=\min_i p_i,\qquad L_{\max}:=\max_i L_i .
\]
\end{assumption}

\section*{Main Convergence Theorem}
\begin{theorem}[Expected suboptimality of RBCD]\label{thm:main}
Under Assumptions \ref{as:convex}–\ref{as:sampling}, run RBCD with a \emph{constant} step size
\[
\eta \le \frac{p_{\min}}{L_{\max}} .
\tag{3}
\]
Then for every $T\ge 1$,
\[
\E\!\bigl[f(w_T)-f(w^\star)\bigr]\;\le\;
\frac{2\,L_{\max}\,R^2}{p_{\min}\,T}.
\tag{4}
\]
Consequently, the method attains an $O(1/T)$ convergence rate in expectation.
\end{theorem}

\section*{Theoretical Properties}
\begin{itemize}[leftmargin=*]
    \item \textbf{Stability.} The bound (4) depends only on the product $L_{\max}R^2$ and the minimal sampling probability $p_{\min}$. Hence the algorithm is robust to the choice of the block partition as long as each block is sampled with non‑vanishing probability.
    \item \textbf{Hyper‑parameter sensitivity.} The step size must satisfy (3). Choosing $\eta$ larger than $p_{\min}/L_{\max}$ can violate the descent inequality (see the proof) and may lead to divergence. The bound is tight up to constant factors for the worst‑case convex quadratic.
    \item \textbf{Comparison to full‑gradient GD.} For $n$ blocks and uniform sampling $p_i=1/n$, (4) becomes
    \[
    \E[f(w_T)-f(w^\star)]\le \frac{2n L_{\max}R^2}{T},
    \]
    which matches the classical $O(1/T)$ rate of gradient descent with step $1/L_{\max}$ up to the factor $n$ reflecting that only one block is updated per iteration.
    \item \textbf{Special case – strongly convex $f$.} If $f$ is $\mu$‑strongly convex, a similar derivation yields a linear (geometric) convergence rate:
    \[
    \E\!\bigl[f(w_T)-f(w^\star)\bigr]\le \bigl(1-\tfrac{\mu p_{\min}}{L_{\max}}\bigr)^{T}\!\bigl[f(w_0)-f(w^\star)\bigr].
    \]
\end{itemize}

\section*{Discussion}
The theorem shows that randomised block updates inherit the optimal $O(1/T)$ rate of full‑gradient methods for convex smooth problems, while each iteration costs only the work needed to evaluate a single block gradient. The proof relies only on convexity and block‑wise smoothness, without any bounded variance or unbiasedness assumptions that are required for stochastic gradient methods.

\newpage
\section*{Preliminaries (Definitions and Lemmas)}
\begin{lemma}[Descent Lemma for a single block]\label{lem:descent}
Let $x\in\R^{d}$ and let block $i$ be selected. For any step $s\in\R^{d_i}$,
\[
f\bigl(x-U_i s\bigr)\le f(x)-\inner{\nabla_i f(x)}{s}+\frac{L_i}{2}\norm{s}^2 .
\tag{5}
\]
\end{lemma}
\begin{proof}
Apply Assumption \ref{as:smooth} with $s\gets -s$ (note $U_i(-s)= -U_i s$). This yields (5). \hfill\qedsymbol
\end{proof}

\section*{Main Proof (Step‑by‑Step Derivation)}
We prove Theorem \ref{thm:main}. Throughout the proof $ \mathcal{F}_t:=\sigma(i_0,\dots ,i_{t-1})$ denotes the sigma‑algebra generated by the past random choices. All expectations are conditional on $\mathcal{F}_t$ unless stated otherwise.

\bigskip
\noindent\textbf{Step 1 (One‑step descent).}  
Condition on the current iterate $w_t$ and the selected block $i_t$. Using Lemma \ref{lem:descent} with $s=\eta \nabla_{i_t} f(w_t)$ we obtain
\[
\begin{aligned}
f(w_{t+1})
&= f\!\bigl(w_t-\eta U_{i_t}\nabla_{i_t}f(w_t)\bigr)                                           \\
&\le f(w_t)-\eta\inner{\nabla_{i_t}f(w_t)}{\nabla_{i_t}f(w_t)}+\frac{L_{i_t}\eta^{2}}{2}
   \norm{\nabla_{i_t}f(w_t)}^{2}                           \\
&= f(w_t)-\eta\Bigl(1-\frac{L_{i_t}\eta}{2}\Bigr)\norm{\nabla_{i_t}f(w_t)}^{2}.
\tag{6}
\end{aligned}
\]
\hfill[Step 1]

\bigskip
\noindent\textbf{Step 2 (Take expectation over the random block).}  
Taking $\E[\;\cdot\mid\mathcal{F}_t]$ on both sides of (6) and using the sampling probabilities $p_i$,
\[
\begin{aligned}
\E\!\bigl[f(w_{t+1})\mid\mathcal{F}_t\bigr]
&\le f(w_t)-\eta\sum_{i=1}^{n}p_i\Bigl(1-\frac{L_i\eta}{2}\Bigr)
   \norm{\nabla_i f(w_t)}^{2}.                                            \\
\end{aligned}
\tag{7}
\]
\hfill[Step 2]

\bigskip
\noindent\textbf{Step 3 (Lower bound on the coefficient).}  
Because $\eta\le p_{\min}/L_{\max}$ (condition (3)) we have for every block $i$,
\[
1-\frac{L_i\eta}{2}\;\ge\;1-\frac{L_{\max}\eta}{2}
\;\ge\;1-\frac{L_{\max}}{2}\cdot\frac{p_{\min}}{L_{\max}}
\;=\;\frac{1+p_{\min}}{2}\;\ge\;\frac{p_{\min}}{2}.
\tag{8}
\]
Thus,
\[
\Bigl(1-\frac{L_i\eta}{2}\Bigr)\ge \frac{p_{\min}}{2},\qquad\forall i .
\]
\hfill[Step 3]

\bigskip
\noindent\textbf{Step 4 (Relate block gradients to full gradient).}  
Recall that $\nabla f(w_t)=\sum_{i=1}^{n}U_i\nabla_i f(w_t)$ and that the blocks are orthogonal, i.e.
$\norm{\nabla f(w_t)}^{2}= \sum_{i=1}^{n}\norm{\nabla_i f(w_t)}^{2}$.
Hence,
\[
\sum_{i=1}^{n}p_i\norm{\nabla_i f(w_t)}^{2}
\;\ge\;p_{\min}\sum_{i=1}^{n}\norm{\nabla_i f(w_t)}^{2}
\;=\;p_{\min}\norm{\nabla f(w_t)}^{2}.
\tag{9}
\]
\hfill[Step 4]

\bigskip
\noindent\textbf{Step 5 (Combine (7), (8) and (9)).}  
Substituting the lower bounds from Steps 3 and 4 into (7) yields
\[
\begin{aligned}
\E\!\bigl[f(w_{t+1})\mid\mathcal{F}_t\bigr]
&\le f(w_t)-\eta\cdot\frac{p_{\min}}{2}\,
       \norm{\nabla f(w_t)}^{2}.
\end{aligned}
\tag{10}
\]
\hfill[Step 5]

\bigskip
\noindent\textbf{Step 6 (Use convexity to replace gradient norm).}  
Convexity (Assumption \ref{as:convex}) implies
\[
f(w_t)-f(w^\star) \;\le\; \inner{\nabla f(w_t)}{w_t-w^\star}
 \;\le\; \norm{\nabla f(w_t)}\;\norm{w_t-w^\star}
 \;\le\; \frac{1}{2\alpha}\norm{\nabla f(w_t)}^{2}
          +\frac{\alpha}{2}\norm{w_t-w^\star}^{2},
\tag{11}
\]
for any $\alpha>0$ (Cauchy–Schwarz + Young’s inequality). Choosing $\alpha = \frac{p_{\min}\eta}{2}$ gives
\[
\frac{p_{\min}\eta}{4}\norm{\nabla f(w_t)}^{2}
\;\ge\; f(w_t)-f(w^\star)-\frac{p_{\min}\eta}{4}\norm{w_t-w^\star}^{2}.
\tag{12}
\]

\bigskip
\noindent\textbf{Step 7 (Insert (12) into (10)).}  
From (10) and (12) we obtain
\[
\begin{aligned}
\E\!\bigl[f(w_{t+1})\mid\mathcal{F}_t\bigr]
&\le f(w_t)-\frac{p_{\min}\eta}{2}\norm{\nabla f(w_t)}^{2}               \\
&\le f(w_t)-2\Bigl(f(w_t)-f(w^\star)\Bigr)
      +\frac{p_{\min}\eta}{2}\norm{w_t-w^\star}^{2}                       \\
&= f(w^\star)+\frac{p_{\min}\eta}{2}\norm{w_t-w^\star}^{2}
      -\bigl(f(w_t)-f(w^\star)\bigr) .
\end{aligned}
\tag{13}
\]
\hfill[Step 7]

\bigskip
\noindent\textbf{Step 8 (Bounding the distance term).}  
From the update rule (1) and the orthogonality of $U_{i_t}$,
\[
\begin{aligned}
\norm{w_{t+1}-w^\star}^{2}
&= \norm{w_t-w^\star -\eta U_{i_t}\nabla_{i_t}f(w_t)}^{2}                \\
&= \norm{w_t-w^\star}^{2}
   -2\eta\inner{U_{i_t}\nabla_{i_t}f(w_t)}{w_t-w^\star}
   +\eta^{2}\norm{U_{i_t}\nabla_{i_t}f(w_t)}^{2}.
\end{aligned}
\tag{14}
\]
Taking expectation conditional on $\mathcal{F}_t$ and using the same arguments as in Steps 2–4,
\[
\E\!\bigl[\norm{w_{t+1}-w^\star}^{2}\mid\mathcal{F}_t\bigr]
\le \norm{w_t-w^\star}^{2}
   -\eta p_{\min}\norm{\nabla f(w_t)}^{2}
   +\eta^{2}L_{\max}p_{\min}\norm{\nabla f(w_t)}^{2}.
\tag{15}
\]
Because $\eta\le p_{\min}/L_{\max}$, the last two terms combine to a non‑positive quantity, implying
\[
\E\!\bigl[\norm{w_{t+1}-w^\star}^{2}\mid\mathcal{F}_t\bigr]
\le \norm{w_t-w^\star}^{2}.
\tag{16}
\]
\hfill[Step 8]

\bigskip
\noindent\textbf{Step 9 (Recursive inequality for the function value).}  
Taking total expectation in (13) and using (16) to drop the distance term yields
\[
\E\!\bigl[f(w_{t+1})\bigr]
\le \E\!\bigl[f(w_t)\bigr]
      -\frac{p_{\min}\eta}{2}\E\!\bigl[\norm{\nabla f(w_t)}^{2}\bigr].
\tag{17}
\]

\bigskip
\noindent\textbf{Step 10 (Relate gradient norm to suboptimality).}  
From convexity and the optimality condition $0\in\partial f(w^\star)$,
\[
f(w_t)-f(w^\star)\;\le\;\inner{\nabla f(w_t)}{w_t-w^\star}
\;\le\;\norm{\nabla f(w_t)}\norm{w_t-w^\star}
\;\le\;\frac{1}{2\eta p_{\min}}\norm{\nabla f(w_t)}^{2}
      +\frac{\eta p_{\min}}{2}\norm{w_t-w^\star}^{2}.
\]
Using (16) to bound $\norm{w_t-w^\star}^{2}\le R^{2}$ and rearranging gives
\[
\frac{p_{\min}\eta}{2}\norm{\nabla f(w_t)}^{2}
\;\ge\; \bigl(f(w_t)-f(w^\star)\bigr)-\frac{p_{\min}\eta R^{2}}{2}.
\tag{18}
\]

\bigskip
\noindent\textbf{Step 11 (Final recursion).}  
Insert (18) into (17):
\[
\E\!\bigl[f(w_{t+1})\bigr]
\le \E\!\bigl[f(w_t)\bigr]
      -\Bigl(\E\!\bigl[f(w_t)-f(w^\star)\bigr]-\frac{p_{\min}\eta R^{2}}{2}\Bigr).
\]
Thus,
\[
\E\!\bigl[f(w_{t+1})-f(w^\star)\bigr]
\le \Bigl(1-\!1\Bigr)\E\!\bigl[f(w_t)-f(w^\star)\bigr]
      +\frac{p_{\min}\eta R^{2}}{2}
= \frac{p_{\min}\eta R^{2}}{2}.
\tag{19}
\]

\bigskip
\noindent\textbf{Step 12 (Summation over $T$ iterations).}  
Summing (19) from $t=0$ to $T-1$ and noting that the left‑hand side telescopes,
\[
\E\!\bigl[f(w_T)-f(w^\star)\bigr]
\le \frac{R^{2}}{2\eta T}\, .
\]
Finally, using the admissible step size bound $\eta = \frac{p_{\min}}{L_{\max}}$ (the largest value satisfying (3)) gives
\[
\E\!\bigl[f(w_T)-f(w^\star)\bigr]
\le \frac{2L_{\max}R^{2}}{p_{\min}T},
\]
which is exactly the claim (4). \hfill\qedsymbol

\section*{Rate Derivation (With Constants)}
The derivation above explicitly shows how each constant appears:

\begin{itemize}
    \item The factor $p_{\min}$ stems from the minimal probability of selecting any block (Step 4).
    \item $L_{\max}$ appears when bounding $1-\frac{L_i\eta}{2}$ from below (Step 3) and when fixing the maximal admissible step size (3).
    \item $R^{2}= \|w_0-w^\star\|^{2}$ is the initial distance (Assumption \ref{as:bounded}) and propagates through the distance recursion (Step 8).
\end{itemize}
Thus the bound (4) is tight up to the universal constant $2$ for the class of convex block‑smooth functions.

\section*{Special Cases}
\begin{enumerate}
    \item \textbf{Uniform sampling.} $p_i=1/n$ for all $i$. Then $p_{\min}=1/n$ and (4) becomes
          \[
          \E[f(w_T)-f(w^\star)]\le \frac{2nL_{\max}R^{2}}{T}.
          \]
    \item \textbf{Strong convexity.} If $f$ is $\mu$‑strongly convex,
          \[
          f(y)\ge f(x)+\inner{\nabla f(x)}{y-x}+\frac{\mu}{2}\norm{y-x}^{2},
          \]
          the same steps lead to the linear rate
          \[
          \E[f(w_T)-f(w^\star)]\le\bigl(1-\tfrac{\mu p_{\min}}{L_{\max}}\bigr)^{T}
               \bigl[f(w_0)-f(w^\star)\bigr].
          \]
    \item \textbf{Non‑uniform probabilities.} Choosing $p_i\propto L_i$ minimizes the constant
          $\frac{L_{\max}}{p_{\min}}$ and yields the optimal bound
          \[
          \E[f(w_T)-f(w^\star)]\le \frac{2\,\bigl(\sum_i L_i\bigr)R^{2}}{T}.
          \]
\end{enumerate}

\end{document}